%% ============================================================
%%  chapters/chapter3.tex - Methodology, Design, and Implementation
%% ============================================================
\chapter{Methodology, Design, and Implementation}
\label{chap:methodology}

\section{Research Methodology}
\label{sec:methodology}

This research uses a design science methodology because the main outcome is an implemented software artifact evaluated against a real educational problem. The process followed five stages:

\begin{enumerate}
  \item \textbf{Problem investigation:} Identify the integrity and learning-process visibility problem in online programming assessments.
  \item \textbf{Literature-informed requirement analysis:} Derive requirements from keystroke dynamics, continuous authentication, programming trace analytics, and assessment governance.
  \item \textbf{Artifact design:} Design a GradeLoop-compatible architecture for keystroke capture, enrollment, verification, monitoring, behaviour analysis, and evidence archiving.
  \item \textbf{Implementation:} Build the module using FastAPI, PyTorch, PostgreSQL, Redis, Docker Compose, and a Next.js evaluator interface.
  \item \textbf{Evaluation:} Evaluate the implementation using real-user BLAIM testing data containing enrollment records, genuine tests, impostor tests, authentication outputs, risk timelines, and archives.
\end{enumerate}

The design science approach is appropriate because the research contribution is not only theoretical. The dissertation demonstrates that an integrated keystroke authentication and behaviour analysis system can be implemented as an operational microservice and assessed using practical metrics.

\section{Functional Requirements}
\label{sec:functional-requirements}

The functional requirements are listed in Table~\ref{tab:functional-requirements}.

\begin{table}[H]
  \centering
  \caption{Functional Requirements of BLAIM}
  \label{tab:functional-requirements}
  \begin{singlespace}
  \begin{tabularx}{\textwidth}{@{}p{2.0cm}p{9.8cm}p{2.2cm}@{}}
    \toprule
    \textbf{ID} & \textbf{Requirement} & \textbf{Status} \\
    \midrule
    FR1 & Capture timestamped keystroke events from programming sessions. & Implemented \\
    FR2 & Enroll students using multi-phase keystroke data. & Implemented \\
    FR3 & Transform raw events into TypeNet-compatible 70 by 5 feature sequences. & Implemented \\
    FR4 & Generate and store 128-dimensional biometric templates. & Implemented \\
    FR5 & Verify a claimed student identity using cosine similarity and a threshold. & Implemented \\
    FR6 & Identify the most likely enrolled user using 1:N matching. & Implemented \\
    FR7 & Monitor active sessions and record timeline risk events. & Implemented \\
    FR8 & Archive completed sessions with events, code, summary risk, and metadata. & Implemented \\
    FR9 & Analyse behaviour for pauses, deletions, paste/copy, burst typing, friction, and engagement. & Implemented \\
    FR10 & Provide instructor/evaluator dashboards for enrollment, testing, metrics, and research review. & Implemented \\
    \bottomrule
  \end{tabularx}
  \end{singlespace}
\end{table}

\section{Non-Functional Requirements}
\label{sec:non-functional-requirements}

The non-functional requirements are:

\begin{itemize}
  \item \textbf{Usability:} Students should type inside a normal coding interface without extra hardware.
  \item \textbf{Performance:} Verification and monitoring should be fast enough for assessment workflows.
  \item \textbf{Scalability:} The service should support class-scale use through stateless API endpoints and external session storage.
  \item \textbf{Reliability:} Session data should be buffered during active work and persisted after submission.
  \item \textbf{Security:} Biometric templates and archives should be protected through database access control, retention limits, and production encryption.
  \item \textbf{Explainability:} The system should expose similarity, risk, confidence, and behaviour indicators rather than opaque decisions only.
  \item \textbf{Maintainability:} The module should follow clear microservice boundaries and documented API contracts.
  \item \textbf{Ethical use:} The system must support informed consent and human review before consequential decisions.
\end{itemize}

\section{System Architecture}
\label{sec:system-architecture}

BLAIM is implemented as a three-layer system: capture, intelligence, and consumption. Figure~\ref{fig:system-architecture} shows the high-level architecture and the communication between the main components.

\begin{figure}[H]
  \centering
  \resizebox{\textwidth}{!}{%
  \begin{tikzpicture}[
    node distance=1.2cm and 1.3cm,
    box/.style={draw, rounded corners, align=center, minimum width=3.0cm, minimum height=0.9cm, fill=gray!8},
    store/.style={draw, cylinder, shape border rotate=90, aspect=0.25, align=center, minimum width=2.2cm, minimum height=1.0cm, fill=blue!8},
    ext/.style={draw, rounded corners, align=center, minimum width=2.8cm, minimum height=0.9cm, fill=green!8},
    arrow/.style={-{Latex[length=2mm]}, thick}
  ]
    \node[box] (ui) {Next.js Coding UI\\Enrollment / Testing};
    \node[box, right=of ui] (eval) {BLAIM Evaluator\\Backend API};
    \node[box, right=of eval] (ks) {Keystroke Service\\FastAPI};
    \node[store, above right=0.8cm and 1.2cm of ks] (redis) {Redis\\Session Buffer};
    \node[store, right=1.4cm of ks] (pg) {PostgreSQL\\Templates / Archives};
    \node[box, below right=0.8cm and 1.2cm of ks] (typenet) {TypeNet Model\\128-d Embeddings};
    \node[ext, above=1.1cm of eval] (dash) {Instructor / Research\\Dashboard};
    \node[ext, below=1.1cm of eval] (rabbit) {Optional RabbitMQ\\Auth Events};
    \node[ext, below=1.1cm of ks] (gemini) {Optional Gemini\\Behaviour Explanation};

    \draw[arrow] (ui) -- node[above, align=center]{keystroke\\batches} (eval);
    \draw[arrow] (eval) -- node[above]{REST} (ks);
    \draw[arrow] (ks) -- (redis);
    \draw[arrow] (ks) -- (pg);
    \draw[arrow] (ks) -- (typenet);
    \draw[arrow] (eval) -- (dash);
    \draw[arrow] (ks) -- (rabbit);
    \draw[arrow] (ks) -- (gemini);
    \draw[arrow] (pg) to[bend left=18] node[above]{records} (dash);

    \node[draw, dashed, fit=(ui), inner sep=0.25cm, label=above:Capture Layer] {};
    \node[draw, dashed, fit=(eval)(ks)(redis)(pg)(typenet)(rabbit)(gemini), inner sep=0.35cm, label=above:Intelligence Layer] {};
    \node[draw, dashed, fit=(dash), inner sep=0.25cm, label=above:Consumption Layer] {};
  \end{tikzpicture}}
  \caption{BLAIM High-Level System Architecture}
  \label{fig:system-architecture}
\end{figure}

\subsection{Capture Layer}
\label{subsec:capture-layer}

The capture layer runs in the browser-based programming interface. It records keystroke events while the student works. Each event includes the student/session identifiers, timestamp, key, key code, dwell time, flight time, assignment identifier, and course identifier when available. The evaluator frontend enforces minimum event counts for enrollment and testing, then sends captured batches to the backend.

\subsection{Intelligence Layer}
\label{subsec:intelligence-layer}

The intelligence layer is the keystroke-service FastAPI application. It performs feature extraction, TypeNet inference, enrollment, verification, identification, continuous monitoring, behaviour analysis, and archive generation. Redis stores active session buffers with a configurable time-to-live. PostgreSQL stores biometric templates, enrollment progress, authentication timeline events, and archives. RabbitMQ publishing is available for authentication result events when enabled.

\subsection{Consumption Layer}
\label{subsec:consumption-layer}

The consumption layer includes evaluator and instructor-facing interfaces. The BLAIM evaluator backend wraps the keystroke-service endpoints and stores local research records. The Next.js evaluator frontend provides enrollment, testing, and research pages that summarize accept/reject outcomes, false positives, false negatives, similarity, risk, and saved metric snapshots.

\subsection{GradeLoop Integration Boundary}
\label{subsec:gradeloop-boundary}

BLAIM is intentionally separated from the core assessment service. The assessment service owns assignment creation, submissions, grading, and course workflows. BLAIM owns keystroke capture processing, biometric templates, session monitoring, and behavioural evidence. This boundary reduces coupling and makes the module reusable in future LMS integrations. The evaluator backend acts as an adapter during research testing: it records local enrollment and test artifacts while forwarding requests to the keystroke service.

\subsection{Data Flow}
\label{subsec:data-flow}

The normal data flow has eight steps:

\begin{enumerate}
  \item The student types in the browser-based coding interface.
  \item The frontend records keystroke timing events and sends batches to the evaluator backend.
  \item The evaluator backend forwards events to the keystroke service.
  \item Redis buffers active session events for monitoring.
  \item The TypeNet pipeline converts recent events into embeddings and computes similarity/risk.
  \item PostgreSQL stores templates, progress, authentication events, and archives.
  \item The evaluator dashboard retrieves metrics, timeline evidence, and per-student records.
  \item Instructors use the outputs for review, intervention, or normal acceptance.
\end{enumerate}

\subsection{External Interfaces}
\label{subsec:external-interfaces}

Table~\ref{tab:external-interfaces} summarizes the communication with external interfaces and infrastructure components.

\begin{table}[H]
  \centering
  \caption{External Interfaces and Communication}
  \label{tab:external-interfaces}
  \begin{singlespace}
  \begin{tabularx}{\textwidth}{@{}p{3.4cm}p{3.5cm}p{5.9cm}@{}}
    \toprule
    \textbf{Interface} & \textbf{Protocol / method} & \textbf{Purpose} \\
    \midrule
    Next.js frontend to evaluator backend & HTTP/JSON & Submit enrollment and testing payloads from the coding UI \\
    Evaluator backend to keystroke service & HTTP/JSON REST & Forward capture, enroll, verify, identify, monitor, finalize, and analytics requests \\
    Keystroke service to Redis & Redis client operations & Buffer active session events and support low-latency monitoring \\
    Keystroke service to PostgreSQL & SQL / JSONB persistence & Store biometric templates, progress rows, timeline events, archives, and metadata \\
    Keystroke service to TypeNet model & Local PyTorch inference & Generate 128-dimensional embeddings from 70 by 5 keystroke sequences \\
    Keystroke service to RabbitMQ & AMQP topic event & Publish authentication result events when event-driven integration is enabled \\
    Keystroke service to Gemini & HTTPS API call & Generate optional natural-language behaviour explanations \\
    Dashboard to evaluator backend & HTTP/JSON & Display enrollment records, tests, snapshots, and research metrics \\
    \bottomrule
  \end{tabularx}
  \end{singlespace}
\end{table}

\section{Keystroke Event Schema}
\label{sec:event-schema}

The core event schema is shown in Table~\ref{tab:event-schema}.

\begin{table}[H]
  \centering
  \caption{Core Keystroke Event Fields}
  \label{tab:event-schema}
  \begin{singlespace}
  \begin{tabularx}{\textwidth}{@{}p{3.0cm}p{2.2cm}p{8.0cm}@{}}
    \toprule
    \textbf{Field} & \textbf{Type} & \textbf{Purpose} \\
    \midrule
    userId & String & Claimed student identifier \\
    sessionId & String & Assessment or enrollment session identifier \\
    timestamp & Integer & Event time in milliseconds \\
    key & String & Key label such as letter, Backspace, Delete, or Enter \\
    keyCode & Integer & Numeric code normalized for model input \\
    dwellTime & Integer & Time the key was held down \\
    flightTime & Integer & Time gap between consecutive key events \\
    assignmentId & String & Optional assessment identifier \\
    courseId & String & Optional course identifier \\
    \bottomrule
  \end{tabularx}
  \end{singlespace}
\end{table}

\section{Feature Extraction and Authentication Pipeline}
\label{sec:feature-pipeline}

Raw keystroke events are transformed into five-feature sequences:

\begin{equation}
\mathbf{f}_i = [HL_i, IL_i, PL_i, RL_i, K_i]
\end{equation}

where $HL$ is hold latency, $IL$ is inter-key latency, $PL$ is press latency, $RL$ is release latency, and $K$ is normalized key code. A sequence contains 70 events:

\begin{equation}
S = [\mathbf{f}_1, \mathbf{f}_2, ..., \mathbf{f}_{70}] \in \mathbb{R}^{70 \times 5}
\end{equation}

The TypeNet-inspired LSTM generates a 128-dimensional embedding:

\begin{equation}
\mathbf{e} = TypeNet(S), \quad \mathbf{e} \in \mathbb{R}^{128}
\end{equation}

During enrollment, the system creates multiple sequences from the captured events. The stored template is the mean embedding:

\begin{equation}
\boldsymbol{\mu}_u = \frac{1}{n}\sum_{i=1}^{n}\mathbf{e}_i
\end{equation}

The template standard deviation is also stored for future calibration:

\begin{equation}
\boldsymbol{\sigma}_u = \sqrt{\frac{1}{n}\sum_{i=1}^{n}(\mathbf{e}_i-\boldsymbol{\mu}_u)^2}
\end{equation}

During verification, the live embedding is compared with the claimed user template using cosine similarity:

\begin{equation}
sim(\mathbf{e}, \boldsymbol{\mu}_u)=
\frac{\mathbf{e}\cdot\boldsymbol{\mu}_u}{\|\mathbf{e}\|\|\boldsymbol{\mu}_u\|}
\end{equation}

The risk score is:

\begin{equation}
risk = 1 - sim
\end{equation}

In the evaluation configuration, the default acceptance threshold is 0.70. Therefore:

\begin{equation}
authenticated =
\begin{cases}
true, & sim \geq 0.70 \\
false, & sim < 0.70
\end{cases}
\end{equation}

\subsection{Sequence Padding and Truncation}
\label{subsec:padding-truncation}

The implemented feature extractor pads shorter sequences with the last available event and truncates long sequences to the most recent 70 events for TypeNet input. This keeps the model input shape stable. For enrollment, the system requires enough captured events to create multiple valid samples, and the authenticator requires at least three valid sequences before a template is created. This protects against incomplete or unrepresentative enrollment samples.

\subsection{Risk Interpretation}
\label{subsec:risk-interpretation}

The risk score is the inverse of similarity. A low risk score indicates that the live embedding is close to the enrolled template. A high risk score indicates a mismatch or model/input problem. Risk is not treated as misconduct by itself. The design deliberately stores the threshold, sample size, session identifier, and confidence so that a reviewer can understand the basis of the decision.

\section{Multi-Phase Enrollment Design}
\label{sec:multi-phase-enrollment}

BLAIM supports phase-based enrollment to reduce dependence on one narrow typing condition. This is necessary because keystroke data is a behavioural biometric rather than a fixed physical trait. The same student can type differently when calm, under time pressure, mentally overloaded, physically tired, unfamiliar with a keyboard, or actively debugging a difficult programming problem. Prior work shows that cognitive and physical stress, emotional state, input device, text length, and task context can measurably change keystroke timing and accuracy \cite{vizer2009stress,epp2011emotional,lee2014emotion,kang2015inputdevices,macdonald2026stress}. In programming specifically, keystroke patterns also change because the task alternates between transcription, planning, syntax recall, debugging, and correction \cite{edwards2020keystroke}.

For this reason, BLAIM does not assume that one enrollment sample represents the whole student. A single calm baseline may produce a clean template, but it may reject the same genuine student during a stressful assessment. A single high-pressure sample may capture time-constrained behaviour, but it may not represent normal coding rhythm. Multi-phase enrollment creates a broader behavioural profile by collecting samples from different but educationally realistic states.

The implemented configuration includes baseline, transcription, stress, and cognitive phases, with baseline and transcription required in the evaluator setup. Table~\ref{tab:enrollment-phases} summarizes the intended purpose of each phase.

\begin{table}[H]
  \centering
  \caption{Enrollment Phases}
  \label{tab:enrollment-phases}
  \begin{singlespace}
  \begin{tabularx}{\textwidth}{@{}p{2.8cm}p{4.2cm}p{6.0cm}@{}}
    \toprule
    \textbf{Phase} & \textbf{Purpose} & \textbf{Example task} \\
    \midrule
    Baseline & Calm-state typing rhythm & Short starter function or normal typing task \\
    Transcription & Motor rhythm with low problem-solving load & Type a fixed Java number utility snippet \\
    Stress & Time-pressure rhythm & Timed FizzBuzz or string reversal challenge \\
    Cognitive & High mental effort typing & Fibonacci memoization or parentheses validation task \\
    \bottomrule
  \end{tabularx}
  \end{singlespace}
\end{table}

The phases have distinct roles. The baseline phase captures the student's normal low-pressure rhythm. The transcription phase controls the content by asking the student to type a fixed code snippet, which helps separate motor typing rhythm from problem-solving difficulty. The stress phase introduces time pressure so that the system can observe how the student's rhythm shifts when tension is present. The cognitive phase captures behaviour during higher reasoning, where long pauses, slower bursts, and correction patterns may be genuine signs of thinking rather than suspicious behaviour.

BLAIM handles situational variation in six ways:

\begin{enumerate}
  \item \textbf{Multiple templates by phase:} Enrollment records are stored with a phase label so that future matching can compare a session against a more relevant condition instead of only one global template.
  \item \textbf{Template variance storage:} Enrollment stores both mean embedding data and template standard deviation. This allows the system to recognize that some students naturally have more variable typing than others.
  \item \textbf{Minimum sample requirements:} The authenticator requires multiple valid sequences before creating a template, reducing the risk that one unusual typing burst becomes the student's identity profile.
  \item \textbf{Rolling monitoring rather than one decision:} Continuous checks are evaluated over session windows, so a temporary pause, syntax struggle, or short stress spike does not automatically define the whole session.
  \item \textbf{Context-aware behaviour metrics:} Pause count, long pauses, deletion rate, burst typing, paste count, rhythm consistency, and friction points are reported beside biometric similarity. This helps reviewers distinguish stress or cognitive load from possible impersonation.
  \item \textbf{Human-in-the-loop interpretation:} BLAIM reports risk and evidence instead of issuing automatic penalties. Device changes, illness, accessibility needs, keyboard problems, stress, or time pressure should be considered during review.
\end{enumerate}

The design goal is therefore not to remove variation from keystroke data. Variation is expected. The goal is to measure it, preserve its context, and use it responsibly so that genuine situational changes do not become unfair authentication failures.

\section{Behaviour Analysis Pipeline}
\label{sec:behaviour-pipeline}

The behaviour analysis engine computes session metrics from the same event stream used for authentication. The main computed metrics are listed in Table~\ref{tab:behaviour-metrics}.

\begin{table}[H]
  \centering
  \caption{Behaviour Metrics Computed from Keystroke Data}
  \label{tab:behaviour-metrics}
  \begin{singlespace}
  \begin{tabularx}{\textwidth}{@{}p{3.3cm}p{4.0cm}p{5.2cm}@{}}
    \toprule
    \textbf{Metric} & \textbf{Computation} & \textbf{Interpretation} \\
    \midrule
    Average typing speed & Typed characters per minute & Fluency and activity level \\
    Deletion rate & Backspace/Delete events divided by total events & Error correction, debugging, or confusion \\
    Pause count & Flight time greater than 1000 ms & Planning, thinking, distraction, or help-seeking \\
    Long pause count & Flight time greater than 3000 ms & Higher cognitive load or inactivity \\
    Paste count & Events labelled as paste & Possible bulk insertion or reuse \\
    Burst typing & Flight time below 100 ms & Very rapid typing, automation, or confident flow \\
    Rhythm consistency & Inverse coefficient of variation of flight time & Stability of typing rhythm \\
    Friction points & 50-event windows with high deletion or long pauses & Localized struggle moments \\
    \bottomrule
  \end{tabularx}
  \end{singlespace}
\end{table}

The behaviour engine computes four instructor-facing scores:

\begin{enumerate}
  \item \textbf{Authenticity score:} Degree to which the session resembles natural human production rather than excessive paste or synthetic activity.
  \item \textbf{Engagement score:} Degree of active participation, based on keystroke volume and timing.
  \item \textbf{Active problem-solving score:} Evidence of incremental construction, editing, and troubleshooting.
  \item \textbf{Learning depth score:} Evidence of meaningful struggle, recovery, and solution development.
\end{enumerate}

The overall process score is:

\begin{equation}
Overall = \frac{Authenticity + Engagement + ProblemSolving + LearningDepth}{4}
\end{equation}

\subsection{Friction Point Detection}
\label{subsec:friction-detection}

BLAIM analyses events in 50-keystroke windows. A window is treated as a friction point when deletion rate is high or long pauses occur repeatedly. The implementation marks friction severity as medium or high based on deletion density. This helps distinguish normal editing from concentrated struggle. In a teaching context, friction points are useful because they identify where the student likely needed to reason, debug, or reconsider the solution.

\subsection{Rule-Based and LLM-Assisted Analysis}
\label{subsec:rule-llm}

The system can operate without an LLM. In that mode, it uses rule-based anomalies and basic pedagogical feedback. If a Gemini API key is configured, the service can generate deeper natural-language explanations from anonymized session metrics and code context. This optional design prevents the core authentication pipeline from depending on external LLM availability while allowing richer feedback in research or low-stakes settings.

\section{Database Design}
\label{sec:database-design}

The keystroke-service PostgreSQL schema contains four main tables:

\begin{itemize}
  \item \textbf{user\_biometrics:} Stores user templates by enrollment phase, including serialized 128-dimensional template data, template standard deviation, sample count, status, and metadata.
  \item \textbf{enrollment\_progress:} Tracks phase completion and overall enrollment status per user.
  \item \textbf{auth\_events:} Stores timeline authentication checks with similarity, risk, confidence, anomaly flags, session offset, threshold, and matched phase.
  \item \textbf{keystroke\_archives:} Stores completed session event logs, final code, summary risk, anomaly count, authentication failures, behavioural analysis, and retention metadata.
\end{itemize}

The evaluator database separately stores enrollment records, test records, metric snapshots, and authentication reruns so that research evidence can be reproduced without rerecording sessions.

\section{API Design}
\label{sec:api-design}

Table~\ref{tab:api-design} lists the main API capabilities.

The API surface is organized around the main lifecycle of a programming assessment session. Capture endpoints receive raw browser events, enrollment endpoints build phase-specific biometric templates, verification and identification endpoints compare live sessions against enrolled profiles, monitoring endpoints support repeated checks during an active session, and finalization endpoints preserve the completed evidence record. This structure keeps the keystroke service usable as a standalone microservice while still allowing the evaluator backend or a future LMS integration to call only the functions required for a particular workflow.

\begin{table}[H]
  \centering
  \caption{Implemented API Capabilities}
  \label{tab:api-design}
  \begin{singlespace}
  \small
  \begin{tabularx}{\textwidth}{@{}p{5.9cm}X@{}}
    \toprule
    \textbf{Endpoint} & \textbf{Purpose} \\
    \midrule
    \path|/api/keystroke/capture| & Buffer keystroke events for an active session \\
    \path|/api/keystroke/enroll/phase| & Enroll a user for a named phase \\
    \path|/api/keystroke/enroll/progress/{user_id}| & Retrieve multi-phase enrollment status \\
    \path|/api/keystroke/verify| & Verify claimed identity using a threshold \\
    \path|/api/keystroke/identify| & Return top matching enrolled users \\
    \path|/api/keystroke/monitor| & Run continuous monitoring over buffered session data \\
    \path|/api/keystroke/session/finalize| & Archive session data after submission \\
    \path|/api/keystroke/timeline/{session_id}| & Retrieve authentication timeline events \\
    \path|/api/keystroke/analytics/{session_id}| & Retrieve risk and behavioural analytics \\
    \path|/ws/monitor/{user_id}/{session_id}| & Stream monitoring updates to dashboards \\
    \bottomrule
  \end{tabularx}
  \end{singlespace}
\end{table}

\section{Technology Choices and Justification}
\label{sec:technology}

Table~\ref{tab:technology} justifies the core technologies used in the implementation.

\begin{table}[H]
  \centering
  \caption{Technology Stack and Justification}
  \label{tab:technology}
  \begin{singlespace}
  \begin{tabularx}{\textwidth}{@{}p{3.0cm}p{3.4cm}p{6.4cm}@{}}
    \toprule
    \textbf{Area} & \textbf{Technology} & \textbf{Justification} \\
    \midrule
    API service & FastAPI & Asynchronous Python API framework with Pydantic validation and OpenAPI support \\
    ML inference & PyTorch & Suitable for LSTM sequence models and loading TypeNet-style weights \\
    Database & PostgreSQL & Reliable relational storage with JSONB support for event archives and metadata \\
    Session buffer & Redis & Low-latency storage for active keystroke sessions with TTL support \\
    Frontend & Next.js/React & Suitable for evaluator pages, coding capture UI, and dashboard interaction \\
    Deployment & Docker Compose & Reproducible orchestration of backend, service, databases, and cache \\
    Messaging & RabbitMQ optional & Allows authentication events to be published to other GradeLoop services \\
    LLM support & Gemini optional & Provides natural-language behavioural explanation when configured \\
    \bottomrule
  \end{tabularx}
  \end{singlespace}
\end{table}

\section{Testing Strategy}
\label{sec:testing-strategy}

The testing strategy combines unit, integration, and scenario-based evaluation:

\begin{itemize}
  \item \textbf{Unit-level validation:} Feature extraction, sequence shaping, template generation, risk calculation, and behaviour metric computation.
  \item \textbf{Integration testing:} End-to-end calls from evaluator frontend to backend to keystroke service and persistence stores.
  \item \textbf{Scenario testing:} Genuine sessions, impostor sessions, low-data sessions, high-risk sessions, and template rebuild cases.
  \item \textbf{Database evidence review:} Verification of enrollment records, test records, archives, auth events, and metric snapshots.
\end{itemize}

\subsection{Evaluation Metrics}
\label{subsec:evaluation-metrics}

The main evaluation metrics are:

\begin{itemize}
  \item \textbf{True accept:} A genuine student session is accepted.
  \item \textbf{False reject:} A genuine student session is rejected.
  \item \textbf{True reject:} An impostor session is rejected.
  \item \textbf{False accept:} An impostor session is accepted.
  \item \textbf{Genuine accept rate:} True accepts divided by genuine tests.
  \item \textbf{False positive rate:} False accepts divided by impostor tests.
  \item \textbf{Average similarity:} Mean biometric similarity over verification outputs.
  \item \textbf{Average risk:} Mean inverse similarity over verification outputs.
\end{itemize}

These metrics were selected because they are understandable to both technical evaluators and academic stakeholders. They also map directly to the risks of the educational domain: incorrectly challenging honest students and incorrectly accepting impostors.

\subsection{Threshold Selection and Calibration}
\label{subsec:threshold-calibration}

The evaluation uses 0.70 as the default acceptance threshold for cosine similarity. This value was selected as a practical prototype threshold rather than a final institutional policy. In the current dataset, it provides a balanced starting point for observing both genuine and impostor outcomes while preserving enough borderline cases for analysis. A threshold that is too low would increase the risk of accepting impostor sessions, while a threshold that is too high would increase false rejection of genuine students whose rhythm changes because of stress, device differences, cognitive load, or task difficulty.

For production deployment, BLAIM should not rely on a single fixed threshold for all students and all assessments. The intended calibration process is:

\begin{enumerate}
  \item collect baseline and transcription enrollment samples for each student;
  \item compute each student's template mean and template variance;
  \item evaluate genuine and impostor validation sessions under the planned assessment condition;
  \item inspect false accept and false reject trade-offs at candidate thresholds;
  \item select institution-level risk bands for low, medium, and high concern;
  \item refine thresholds per phase or per student when enough historical data is available.
\end{enumerate}

This makes the 0.70 threshold an evaluation configuration for the prototype, while the long-term design supports adaptive calibration. The reviewer should therefore interpret the reported accuracy metrics as feasibility evidence, not as a final deployment guarantee.

\subsection{Requirement Traceability}
\label{subsec:requirement-traceability}

Table~\ref{tab:traceability} links key requirements to implementation artifacts and evaluation evidence.

\begin{table}[H]
  \centering
  \caption{Requirement Traceability Matrix}
  \label{tab:traceability}
  \begin{singlespace}
  \scriptsize
  \begin{tabularx}{\textwidth}{@{}p{1.2cm}p{3.8cm}p{4.4cm}p{3.2cm}p{1.4cm}@{}}
    \toprule
    \textbf{ID} & \textbf{Requirement} & \textbf{Implementation artifact} & \textbf{Evaluation evidence} & \textbf{Status} \\
    \midrule
    FR1 & Capture keystroke events & Next.js capture UI and \path|/api/keystroke/capture| & 31 test records with 105--421 events & Passed \\
    FR2 & Multi-phase enrollment & \path|/api/keystroke/enroll/phase|, \path|enrollment_progress| & 32 enrollment records, 15 progress rows & Passed \\
    FR3 & Feature transformation & \texttt{create\_typenet\_sequence} feature extractor & 70 by 5 model input pipeline & Passed \\
    FR4 & Biometric template storage & \path|user_biometrics| table & 32 stored biometric templates & Passed \\
    FR5 & Claimed-user verification & \path|/api/keystroke/verify| & 31 verification responses & Passed \\
    FR6 & Top-K identification & \path|/api/keystroke/identify| & Identify responses stored in evaluator records & Passed \\
    FR7 & Continuous monitoring & \path|/api/keystroke/monitor|, timeline view & 12 authentication timeline events & Passed \\
    FR8 & Session archiving & \path|/api/keystroke/session/finalize| & 34 keystroke archives & Passed \\
    FR9 & Behaviour analysis & \path|/api/keystroke/analyze|, analytics endpoint & Behaviour metrics engine implemented & Passed \\
    FR10 & Research dashboard & Next.js research page, evaluator backend & 2 saved metric snapshots & Passed \\
    \bottomrule
  \end{tabularx}
  \end{singlespace}
\end{table}

\section{Commercialization Plan}
\label{sec:commercialization}

BLAIM has commercialization potential as an academic integrity and learning analytics add-on for programming assessment platforms. The value proposition is a privacy-conscious alternative or complement to camera proctoring: continuous authorship confidence and process evidence collected directly from the coding environment.

Potential customer groups include universities, online degree providers, coding bootcamps, certification providers, and LMS vendors. A feasible commercialization model is software-as-a-service with institution-level licensing, plus an on-premise option for institutions with strict data governance policies.

The product would be positioned as a specialist module for programming assessments rather than a general proctoring tool. This positioning is important because many existing academic-integrity products focus on identity checks, lockdown browsers, webcam monitoring, or final-submission similarity. BLAIM's commercial value is the process layer: it shows how code was produced while also maintaining continuous authorship confidence.

The differentiators are:

\begin{itemize}
  \item Continuous authentication without webcam dependence.
  \item Behaviour analysis that supports learning intervention as well as integrity review.
  \item API-first microservice design for LMS and coding-platform integration.
  \item Evidence archives for instructor review and appeals.
  \item Configurable thresholds and retention policies for institutional governance.
\end{itemize}

Commercial readiness requires additional work: production encryption of biometric templates, fairness audits, administrator policy controls, consent workflows, and formal security testing. The minimum viable product would package the keystroke service, evaluator dashboard, LMS integration adapter, policy configuration panel, and evidence export workflow. A pilot deployment could begin in a single programming module, then expand after threshold calibration, instructor training, and student consent processes are validated.

Table~\ref{tab:commercial-positioning} compares BLAIM against common assessment-integrity approaches.

\begin{table}[H]
  \centering
  \caption{Commercial Positioning Against Existing Integrity Approaches}
  \label{tab:commercial-positioning}
  \begin{singlespace}
  \scriptsize
  \begin{tabularx}{\textwidth}{@{}p{2.7cm}p{3.0cm}p{3.8cm}p{3.7cm}@{}}
    \toprule
    \textbf{Approach} & \textbf{Primary evidence} & \textbf{Limitation in coding assessments} & \textbf{BLAIM advantage} \\
    \midrule
    Webcam proctoring & Student and room observation & Intrusive and weak at explaining how code was produced & Less visually intrusive; captures authorship process inside the coding interface \\
    Lockdown browser & Environment restriction & Cannot prove who typed after login or how code evolved & Adds continuous identity and session-process evidence \\
    Code similarity checking & Final submitted code & Misses private, generated, rewritten, or heavily modified sources & Shows whether code appeared through natural typing, debugging, or paste-like activity \\
    Login authentication & Initial identity check & Does not detect mid-session typist changes & Performs repeated verification and monitoring during the session \\
    Learning analytics only & Process and engagement traces & Usually lacks biometric identity continuity & Combines learning-process indicators with identity confidence \\
    \bottomrule
  \end{tabularx}
  \end{singlespace}
\end{table}

The proposed revenue model has two tiers. A SaaS tier would suit institutions that want low operational overhead and automatic updates. An institution-hosted tier would suit universities with strict privacy, data residency, or biometric governance requirements. Commercial adoption should include implementation services for LMS integration, threshold calibration, staff training, and policy configuration.

Table~\ref{tab:commercialization-evidence} summarizes the commercialization evidence.

\begin{table}[H]
  \centering
  \caption{Commercialization Evidence}
  \label{tab:commercialization-evidence}
  \begin{singlespace}
  \begin{tabularx}{\textwidth}{@{}p{3.0cm}p{4.4cm}p{5.6cm}@{}}
    \toprule
    \textbf{Dimension} & \textbf{Evidence} & \textbf{Commercial value} \\
    \midrule
    Customer problem & Online programming assessments need identity assurance without relying only on webcam surveillance & Addresses a current integrity challenge in blended and remote education \\
    Product capability & Continuous keystroke verification, behaviour analysis, archives, and dashboards & Provides both academic integrity evidence and learning-process insight \\
    Differentiation & Works inside the coding environment and produces process evidence & Less intrusive than camera-only proctoring and richer than final-code similarity checks \\
    Integration path & REST APIs, WebSocket monitoring, PostgreSQL, Redis, optional RabbitMQ & Can be integrated with LMSs, coding platforms, and institutional dashboards \\
    Revenue model & SaaS subscription or institution-hosted deployment & Supports both cloud-first and privacy-sensitive institutions \\
    Expansion path & LMS plugins, adaptive thresholds, multimodal telemetry, and analytics reports & Provides a roadmap beyond the research prototype \\
    \bottomrule
  \end{tabularx}
  \end{singlespace}
\end{table}

\clearpage
